
%%%%%%%%%%%%%%%%%%%%%%%%%%%%%%%%%%%%%%%%%%%%%%%%%%%%%%%%%%%%%%
%%%%%% MA-0471 Introducción a la Geometría Diferencial %%%%%%%
%%%%%%%%%%%%%%%%%%%%%%%%%%%%%%%%%%%%%%%%%%%%%%%%%%%%%%%%%%%%%%
\newpage
\subsection*{MA-0471 Introducción a la Geometría Diferencial}
\addcontentsline{toc}{subsection}{MA-0471 Introducción a la Geometría Diferencial}

\hypertarget{MA0471}{}
\begin{refsection}
\begin{table}[H]
\centering{
\begin{tabular}{|ll|}
\hline
%\multicolumn{2}{|c|}{\textbf{Nivel de virtualidad del curso:} Virtual}\\ \hline
\textbf{Año:} II  & \textbf{Requisitos:} MA-0351  \\ 
\textbf{Ciclo:} IV  & \textbf{Correquisitos:} Ninguno \\ 
\textbf{Tipo de curso:} Teórico & \textbf{Horas presenciales por semana:}   \\ 
\textbf{Créditos:} 3 & \textbf{Horas de trabajo independiente por semana:}  \\ \hline
\end{tabular}
}
\end{table}



\subsubsection*{Descripción del curso}

El presente curso ofrece un estudio riguroso e intuitivo de las propiedades geométricas de curvas y superficies en el espacio, haciendo uso de herramientas analíticas. A lo largo del curso, se desarrollan los fundamentos necesarios para comprender la geometría local de curvas mediante su parametrización, curvatura y torsión, así como el comportamiento global de las mismas. Posteriormente, se introduce el estudio de superficies, abordando su descripción paramétrica, regularidad, plano tangente, área y propiedades topológicas como la orientabilidad y la compacidad.

Asimismo, el curso profundiza en la geometría extrínseca a través de la aplicación de Gauss y la curvatura gaussiana, y en la geometría intrínseca mediante el análisis de isometrías, el Teorema Egregium, el transporte paralelo y las geodésicas.  Este es un tema mucho m\'as profundo que el de curvas, tanto a nivel filos\'ofico como a nivel matem\'atico. El enfoque seguido en el curso es el Riemanniano, es decir, el que propuso Bernhard Riemann es su lectura inaugural en la Universidad de Gotingen: \textit{ On the Hypotheses Which Lie at the Foundations of Geometry  }. 

El curso pretende también servir como una introducción concreta al curso más avanzado y abstracto de Geometría Diferencial que estudia las variedades diferenciales con mayor generalidad. 


Se promueve un enfoque que equilibra la formalidad matemática con la intuición geométrica, incorporando ejemplos, visualizaciones y uso de herramientas tecnológicas. Además, se fomenta el desarrollo de habilidades en demostración matemática, pensamiento crítico y trabajo colaborativo, preparando al estudiantado para cursos más avanzados y aplicaciones en diversas áreas de la matemática.


\subsubsection*{Objetivos}

\textbf{General:} Aplicar de forma rigurosa pero intuitiva la geometría diferencial de curvas y superficies, mediante el estudio de sus propiedades locales y globales, tanto extrínsecas como intrínsecas. \\

\textbf{Específicos:}
Al finalizar el curso se espera que la persona estudiante sea capaz de:
\begin{enumerate}
\item Analizar la geometría de curvas en el espacio mediante su parametrización, regularidad, longitud de arco, y el uso del triedro de Frenet, curvatura y torsión.


\item Describir superficies a través de parametrizaciones locales, condiciones de regularidad, plano tangente y herramientas como el teorema del valor regular.


\item Demostrar resultados involucrando las propiedades geométricas y topológicas de las superficies, incluyendo área, orientabilidad, compactación e intersecciones.


\item Aplicar la teoría de la aplicación de Gauss, la segunda forma fundamental y la curvatura gaussiana para caracterizar la geometría extrínseca de las superficies, incluyendo casos especiales como superficies minimales y regladas.


\item Analizar la geometría intrínseca de superficies mediante isometrías, aplicaciones conformes, el Teorema Egregium, y el estudio de transporte paralelo y geodésicas.
\end{enumerate}

\subsubsection*{Contenidos}

\begin{enumerate}
\item \textbf{Curvas (3 semanas):} 
\begin{enumerate}
    \item Parametrización, regularidad, longitud de arco. 
    \item Triedro móvil, curvatura y torsion. 
    \item Fórmulas de Frenet. Forma canónica. 
    \item Propiedades globales: el teorema de los 4 vértices.
\end{enumerate}

\item \textbf{Superficies (3.5 semanas):} 
\begin{enumerate}
    \item Regularidad y parametrización local. 
    \item Teorema del valor regular. 
    \item Funciones derivables en una superficie. Plano tangente. 
    \item Inmersiones y submersiones, intersección transversal para superficies. 
    \item Forma de área. Definición geométrica de área.
    \item Orientabilidad. 
    \item Caracterización de superficies compactas orientables. 
\end{enumerate}

\item \textbf{Aplicación de Gauss (3.5 semanas):} 
\begin{enumerate}
    \item Definición y propiedades. 
    \item Segunda forma fundamental. 
    \item Curvatura Gaussiana. 
    \item Campos vectoriales en superficies. Campos vectoriales como operadores
diferenciales de orden 1. 
    \item Superficies regladas y superficies minimales.
\end{enumerate}

\item \textbf{Geometría intrínseca de superficies: (3 semanas):} 
\begin{enumerate}
    \item Isometrías y aplicaciones conformes. 
    \item Teorema Egregium. 
    \item Transporte paralelo y geodésicas
\end{enumerate}

\item \textbf{Geometría proyectiva (2 semanas):} 
\begin{enumerate}
    \item 
\end{enumerate}


\end{enumerate}

\subsubsection*{Metodología}

\noindent \textbf{Lineamientos metodológicos:} La persona docente a cargo de este curso debe respetar las siguientes pautas:
\begin{enumerate}
    \item La presentación de los contenidos debe adaptarse a un contexto donde los estudiantes tienen habilidades computacionales avanzadas pero están avanzando en su preparación en demostraciones rigurosas. 
    \item Algunos aspectos del curso requieren, para su total rigurosidad, del resultado fundamental de existencia y unicidad para ecuaciones diferenciales. Este resultado no ha sido probado en segundo año, por lo que el mismo debe enunciarse y explicarse, pero no probarse, en este curso.
    \item En este curso se aplican las bases de lógica y las estrategias de demostración vistas en el curso MA-0252 Introducción a las demostraciones.
    \item Las demostraciones deben presentarse tomando en cuenta que están dirigidas a una persona estudiante de formación intermedia. En particular, se debe establecer rigurosidad pero mantener un contexto intuitivo para los contenidos, 
    \item Deben incluirse, cuando el material lo permita, abundantes ejemplos para ayudar a la comprensión.
    \item Cuando la naturaleza del contenido lo permita, se deben utilizar representaciones visuales que apoyen la comprensión de los temas.
    \item La persona docente debe tomar en cuenta que el nivel de dedicación del estudiantado al curso debe ser compatible con el creditaje del mismo.
\end{enumerate}

\textbf{Sugerencias metodológicas:} Adicionalmente, se tienen las siguientes recomendaciones para la persona docente a cargo de este curso:
\begin{enumerate}
    \item Utilizar el recurso tecnológico para reforzar distintos conceptos estudiados en el curso por medio de la visualización, cálculo y experimentación.
    \item Aprovechar momentos adecuados del curso para motivar el estudio por medio de reseñas acerca del desarrollo histórico del análisis matemático y su importancia en aplicaciones dentro y fuera de la matemática
    \item En la evaluación, se sugiere utilizar estrategias que promuevan el trabajo constante de parte de las personas estudiantes a lo largo del ciclo lectivo. Por ejemplo, esto se puede hacer por medio de tareas o quizzes.
    \item Para fomentar el desarrollo de las habilidades investigativas en el estudiantado, se sugiere la implementación de pequeños proyectos de investigación. Por ejemplo, pueden ser proyectos de investigación bibliográfica o también proyectos cortos donde se desarrolle un tema por medio de ejercicios dirigidos.
    \item Dado que hoy en día la mayoría del trabajo en matemática, tanto a nivel académico como profesional, es de naturaleza colaborativa, se sugiere a la persona docente implementar estrategias que promuevan el trabajo en equipo.
\end{enumerate}



%%%%%%%%%%%%%%%%%%%%%%%%%%%%%%%%%%
%%%%%%%%%%%%%%%%%%%%%%%%%%%%%%%%%%
%% No cite para las referencias %%
\nocite{DoC16}
\nocite{Kuh15}
\nocite{ONeil06}
%%%%%%%%%%%%%%%%%%%%%%%%%%%%%%%%%%
%%%%%%%%%%%%%%%%%%%%%%%%%%%%%%%%%%


\printbibliography[heading=subbibliography]
\end{refsection}
